% Part: proof-theory
% Chapter: natural-deduction
% Section: rules-proofs

\documentclass[../../../include/open-logic-section]{subfiles}

\begin{document}

\olfileid{pt}{nat}{rp}

\olsection{नियम आणि \usetoken{P}{proof}}

प्रथम काही व्याख्या करू
आणि संज्ञा निश्चित करू.

~\Log{N1c} आणि~\Log{N1i}
मधील $!A \lif !B$
उल्लेखणारे दोन नियम पाहू:
\[
\bottomAlignProof
\AxiomC{$\Discharge{!A}{x}$}
\shortDeduce
\DeduceC{$!B$}
\DischargeRule{\Intro{\lif}}{x}
\UnaryInfC{$!A \lif !B$}
\DisplayProof
\qquad
\bottomAlignProof
\AxiomC{$!A \lif !B$}
\AxiomC{$!A$}
\RightLabel{\Elim{\lif}}
\BinaryInfC{$!B$}
\DisplayProof
\]
\Elim{\lif} नियम सोपा आहे.
रेषेच्या वरची !!{formula}s
ही \emph{आधारविधाने} होत.
डावीकडील !!{formula}
हे \emph{मुख्य}
!!{formula}~$!A \lif !B$
आहे. सर्व !!{elimination}
नियमांत असे एक आधारविधान
असते आणि ते नेहमी
सर्वात डावीकडे असते.
त्याला अनुमानाचे
\emph{प्रमुख} आधारविधान
म्हणतात. या उदाहरणाप्रमाणे
इतर आधारविधाने असतील,
तर त्यांना \emph{गौण}
आधारविधाने म्हणतात.
!!{introduction} नियमांची
आधारविधानेही प्रमुख
आधारविधाने म्हणतात.
रेषेखालील !!{formula}
हा \emph{निष्कर्ष} आहे.

\Intro{\lif} नियमाला
एकच आधारविधान असते;
येथे निष्कर्ष हे
मुख्य सूत्र~$!A \lif !B$
असते. सर्व !!{introduction}
नियम असेच असतात:
निष्कर्ष हे मुख्य सूत्र
असते, आणि त्याचा
मुख्य संकारकच “प्रवेशित”
केला जातो.

~$\Discharge{!A}{x}$
या लेखनाचा अर्थ पुढीलप्रमाणे.
एखाद्या !!a{proof} मध्ये
\Intro{\lif} अनुमान
आधारविधान म्हणून
!!a{formula}~$!B$
वर लावतात.
~$!B$ मध्ये संपणाऱ्या
उप-!!{proof} च्या
सुरुवातीला काही
!!{formula}s असतात;
त्यांना \emph{गृहीतके}
म्हणतात. ~$!A$
रूपाची एक किंवा अधिक
गृहीतके धरून $!B$
ची !!^a{proof} म्हणजे
$!A$ वरून~$!B$
निष्पन्न होते हे दाखवणारी
!!a{proof} होय.
\Intro{\lif} लावल्यावर
$!A \lif !B$
निष्पन्न होते; आता
निष्कर्ष गृहीतक~$!A$
वर अवलंबून राहत नाही.
हे दाखवण्यासाठी,
\Intro{\lif} असलेल्या
!!{proof} मधील
$!A$ रूपाची संबंधित
गृहीतके \emph{मुक्त}
किंवा \emph{बंद} केली
जातात. कोणती गृहीतके
मुक्त करता येतात, हे
नियम सांगतो.
$!A \lif !B$ निष्कर्ष
असलेल्या \Intro{\lif}
मध्ये ती $!A$ रूपाची,
म्हणजे सोपाधिक विधानाच्या
पूर्वांगाशी जुळणारी
गृहीतके असतात.
कोणत्या नियमाने कोणती
गृहीतके मुक्त होतात हे
दाखवणे कधी महत्त्वाचे
असते; त्यासाठी मुक्तीची
खूण~$x$ वापरतात.
मात्र गृहीतके मुक्त
करणाऱ्या नियमाने योग्य
रूपाची \emph{सर्वच}
गृहीतके, किंवा एकही
गृहीतक, मुक्त करणे
\emph{बंधनकारक नाही}.
उदाहरणार्थ, पहिल्या
\Intro{\lif} ने एकही
गृहीतक मुक्त केलेले
नसले तरी पुढील सिद्धता
योग्य आहे:
\[
\AxiomC{$\Discharge{!A}{y}$}
\DischargeRule{\Intro{\lif}}{x}
\UnaryInfC{$!B \lif !A$}
\DischargeRule{\Intro{\lif}}{y}
\UnaryInfC{$!A \lif (!B \lif !A)$}
\DisplayProof
\]
नियमाने एकही गृहीतक
मुक्त केले नाही, तर
मुक्तीची खूण वगळतो
(येथेही तसे केले आहे).

पुढील !!{proof} पाहा:
\[
\AxiomC{$\Discharge{!A}{x}$}
\AxiomC{$\Discharge{!A}{y}$}
\RightLabel{\Intro{\land}}
\BinaryInfC{$!A \land !A$}
\RightLabel{\Elim{\land}}
\UnaryInfC{$!A$}
\DischargeRule{\Intro{\lif}}{x}
\UnaryInfC{$!A \lif !A$}
\DischargeRule{\Intro{\lif}}{y}
\UnaryInfC{$!A \lif (!A \lif !A)$}
\DisplayProof
\]
पहिले \Intro{\lif}
अनुमान डावीकडील
गृहीतक~$!A^x$ मुक्त करते;
दुसरे उजवीकडील
गृहीतक~$!A^y$ मुक्त करते.
म्हणजे खूण~$x$ असलेली
सर्व गृहीतके खूण~$x$
असलेल्या \Intro{\lif}
ने, आणि खूण~$y$
असलेली सर्व गृहीतके
खूण~$y$ असलेल्या
त्या अनुमानाने
मुक्त होतात.
यासाठी एकाच खूणेची
सर्व गृहीतके तेच
!!{formula} असणे
महत्त्वाचे आहे.

संख्यापकांचे नियम
किंचित अधिक गुंतागुंतीचे आहेत.
उदाहरणार्थ,~\Log{N1i}
मध्ये $\lexists$
साठीचे नियम असे:
\[
\bottomAlignProof
\AxiomC{$!A(t)$}
\RightLabel{\Intro{\lexists}}
\UnaryInfC{$\lexists[x][!A(x)]$}
\DisplayProof
\qquad
\bottomAlignProof
\AxiomC{$\lexists[x][!A(x)]$}
\AxiomC{$\Discharge{!A(c)}{x}$}
\shortDeduce
\DeduceC{$!B$}
\DischargeRule{\Elim{\lexists}}{x}
\BinaryInfC{$!B$} \DisplayProof
\bottomAlignProof
\]
\Intro{\lexists} मध्ये
आधारविधान~$!A(t)$ असते;
येथे $t$ हे बंद पद,
म्हणजे चल नसलेले पद आहे.
असे अनुमान योग्य,
म्हणजे \Intro{\lexists}
या नियम-रूपबंधाचे उदाहरण,
तेव्हाच असते जेव्हा
एखादे !!{formula}~$!A$,
चल~$x$ आणि बंद पद~$t$
असते ज्यामुळे निष्कर्ष
$\lexists[x][!A]$
आणि आधारविधान
~$\Subst{!A}{t}{x}$
असते. \Elim{\lexists}
नियमाने $!A(c)$
रूपाची गृहीतके मुक्त
करता येतात: प्रत्येक
अनुमानासाठी असा एखादा
!!{constant}~$c$ असतो
की $\Subst{!A}{c}{x}$
रूपाची गृहीतके मुक्त
होतात. एका अनुमानाने
मुक्त होणाऱ्या सर्व
गृहीतकांत तोच~$c$
वापरला पाहिजे.
या स्थिरांकाला
\Elim{\lexists}
अनुमानाचे
\emph{आयगेन चर} म्हणतात.
~$c$ हे मुक्त होणाऱ्या
~$!A(c)$ व्यतिरिक्त
इतर उघड्या गृहीतकात,
~$!B$ मध्ये किंवा
~$!A$ मध्ये येत
नसेल तेव्हाच अनुमान
योग्य असते. याला
\emph{आयगेन चराची अट}
म्हणतात.

~\Log{N1c} आणि~\Log{N1i}
मधील !!^{proof}s म्हणजे
नियम वापरून स्वयंसिद्धकांपासून
विगमनाने तयार होणारे
सूत्रांचे सांत वृक्ष.
प्रत्येक पान हे मुक्तीच्या
खुणेसह किंवा तिच्याविना
असलेले गृहीतक असते;
इतर प्रत्येक !!{formula}
त्याच्या लगत वरच्या
!!{formula}s पासून
पद्धतीच्या अनुमान नियमाने
निष्पन्न होते.
अनुमानाच्या वरच्या
वृक्षातील गृहीतक म्हणून
आलेले एखादे !!{formula}
त्या अनुमानापर्यंत मुक्त
झालेले नसेल, तर त्याला
(त्या अनुमानाशी संबंधित)
\emph{उघडे} म्हणतात.

\begin{defn}\ollabel{defn:proof}
~\Log{N1c} आणि~\Log{N1i}
मधील \emph{!!^{proof}s}
हे पुढीलप्रमाणे विगमनाने
परिभाषित !!{formula}s
चे सांत वृक्ष आहेत:
\begin{enumerate}
  \item\ollabel{defn:prf:assumption}
  मुक्तीची $x$ किंवा $y$
  सारखी खूण असलेले
  प्रत्येक !!{formula}
  स्वतःच !!a{proof} असते.
  एकाच !!{formula}
  ची !!a{proof} असेल,
  तर ते !!{formula}
  उघडे गृहीतक मानतात.
  \item\ollabel{defn:prf:one-prem}
  $R$ हा एक, दोन किंवा
  तीन आधारविधाने
  $!A_1$, $!A_2$, $!A_3$
  आणि निष्कर्ष~$!A$
  असलेला नियम असू द्या.
  $\delta_1$, $\delta_2$,
  $\delta_3$ या अनुक्रमे
  $!A_1$, $!A_2$,
  $!A_3$ मध्ये संपणाऱ्या
  !!{proof}s असतील, तर
  पुढीलपैकी योग्य रचना
  अनुक्रमे !!a{proof} असते:
  \[
  \AxiomC{}
  \RightLabel{$\delta_1$}
  \DeduceC{$!A_1$}
  \RightLabel{$R$}
  \UnaryInfC{$!A$}
  \DisplayProof
\qquad
  \AxiomC{}
  \RightLabel{$\delta_1$}
  \DeduceC{$!A_1$}
  \AxiomC{}
  \RightLabel{$\delta_2$}
  \DeduceC{$!A_2$}
  \RightLabel{$R$}
  \BinaryInfC{$!A$}
  \DisplayProof
\qquad
  \AxiomC{}
  \RightLabel{$\delta_1$}
  \DeduceC{$!A_1$}
  \AxiomC{}
  \RightLabel{$\delta_2$}
  \DeduceC{$!A_2$}
  \AxiomC{}
  \RightLabel{$\delta_3$}
  \DeduceC{$!A_3$}
  \RightLabel{$R$}
  \TrinaryInfC{$!A$}
  \DisplayProof
  \]
  मात्र त्या !!{proof}s
  मधील गृहीतक-खुणा
  सुसंगत असल्या पाहिजेत
  (दोन भिन्न
  !!{formula}s ना
  तीच मुक्तीची खूण
  नसावी), आणि ही
  अनुमाने~$R$ चे
  योग्य उपयोग असले पाहिजेत.
  \item नव्या !!{proof}
  मधील सर्वात वरची
  !!{formula} उदाहरणे
  ही त्या !!{proof}
  ची \emph{उघडी गृहीतके}
  मानली जातात, जर ती
  $\delta_1$, $\delta_2$
  किंवा~$\delta_3$
  मधील उघडी गृहीतके
  असतील आणि नियमाने
  मुक्त केलेली नसतील.
  नियमावर खूण~$x$
  (किंवा $x\, y$)
  असेल, तर \Intro{\lif}
  आणि~\Intro{\lnot}
  साठी $\delta_1$
  मधील खूण~$x$
  असलेली सर्व उघडी
  गृहीतके मुक्त होतात;
  \Elim{\lexists}
  साठी $\delta_2$
  मधील खूण~$x$
  असलेली सर्व उघडी
  गृहीतके मुक्त होतात;
  \Elim{\lor} साठी
  $\delta_2$ मधील
  खूण~$x$ असलेली
  आणि~$\delta_3$ मधील
  खूण~$y$ असलेली
  सर्व उघडी गृहीतके
  मुक्त होतात.
\end{enumerate}
\end{defn}

!!a{proof} ची
विगमनात्मक रचना करताना
वापरलेल्या !!{proof}s ना
तिच्या \emph{उप-!!{proof}s}
म्हणतात. अनुमानाच्या
आधारविधानांत संपणाऱ्या
उप-!!{proof}s ना त्या
अनुमानाच्या \emph{निकटच्या
उप-!!{proof}s} म्हणतात.
\Log{N1c} आणि \Log{N1i}
यांचे नियम
\olref{tab:N1} मध्ये आहेत.

\subfile{rules-N1}

\begin{defn}
!!a{proof}~$\delta$ ची
\emph{!!{height}}~
$\pheight{\delta}$
पुढीलप्रमाणे विगमनाने
परिभाषित करतात.
\begin{enumerate}
  \item $\delta$ केवळ
  गृहीतकाची असेल, तर
  $\pheight{\delta} = 0$.
  \item $\delta$ चा शेवट
  एका आधारविधानाच्या
  अनुमानात होत असेल
  आणि निकटची
  उप-!!{proof}~$\delta_1$
  असेल, तर
  $\pheight{\delta} =
  \pheight{\delta_1} + 1$.
  \item $\delta$ चा शेवट
  दोन आधारविधानांच्या
  अनुमानात होत असेल
  आणि निकटच्या
  उप-!!{proof}s~$\delta_1$
  आणि~$\delta_2$ असतील,
  तर $\pheight{\delta} =
  \max(\pheight{\delta_1},
  \pheight{\delta_2}) + 1$.
  \item $\delta$ चा शेवट
  \Elim{\lor} मध्ये होत
  असेल आणि निकटच्या
  उप-!!{proof}s~$\delta_1$,
  $\delta_2$ आणि~$\delta_3$
  असतील, तर
  $\pheight{\delta} =
  \max(\pheight{\delta_1},
  \pheight{\delta_2},
  \pheight{\delta_3}) + 1$.
\end{enumerate}
क्रमवर्ती~$S$ ची उंची
$\le n$ असलेली
!!a{proof} असेल,
तर $ \Proves[n] S$
असे लिहितो.
\end{defn}

\begin{ex}
~\Log{N1i} मध्ये कोणतेही
!!{formula} हे स्वतःपासून
उघडे गृहीतक घेऊन
!!a{proof} असते.
म्हणून
\[
!A \land !B^x
\]
ही !!a{proof}~$\delta_1$
आहे; तिची उंची~$0$.
~$\delta_1$ वर
~\Elim{\land} च्या
दोनपैकी एक आवृत्ती
लावून !!{proof}s~$\delta_2$
आणि~$\delta_3$ मिळतात:
\[
\AxiomC{$!A \land !B^x$}
\RightLabel{\Elim{\land}[2]}
\UnaryInfC{$!B$}
\DisplayProof
\qquad
\AxiomC{$!A \land !B^x$}
\RightLabel{\Elim{\land}[1]}
\UnaryInfC{$!A$}
\DisplayProof
\]
~$\delta_2$ आणि
~$\delta_3$ मधील
अंतिम अनुमानांच्या निकटच्या
उप-!!{proof}s सारख्याच:
$!A \land !B$ या
गृहीतकापासून तयार झालेल्या.
दोन्ही !!{proof}s मध्ये
$!A\land !B$ हे
उघडे गृहीतक म्हणून येते.
आपल्याला
$\pheight{\delta_2} =
\pheight{\delta_3} = 1$
मिळते.

~\Intro{\land} नियमात
$!B$ आणि $!A$
या रूपाची दोन
आधारविधाने असतील,
तर $!B \land !A$
या रूपाचा निष्कर्ष
निघतो. त्यामुळे
पुढील !!a{proof}~$\delta_4$
आहे:
\[
\AxiomC{$!A \land !B^x$}
\RightLabel{\Elim{\land}[2]}
\UnaryInfC{$!B$}
\LeftSubproofLabel{$\delta_2$}
\AxiomC{$!A \land !B^x$}
\RightLabel{\Elim{\land}[1]}
\UnaryInfC{$!A$}
\RightSubproofLabel{$\delta_3$}
\RightLabel{\Intro{\land}}
\BinaryInfC{$!B \land !A$}
\DisplayProof
\]
तिची उंची
$\pheight{\delta_4} =
\max(\pheight{\delta_2},
\pheight{\delta_3})+1
= \max(1,1) + 1 = 2$.
तिच्यात $!A \land !B$
या गृहीतकाची दोन
उदाहरणे आहेत; दोन्ही
उघडी आहेत.

शेवटी~\Intro{\lif}
लावून $\delta_5$
मिळवू:
\[
\AxiomC{$\Discharge{!A \land !B}{x}$}
\RightLabel{\Elim{\land}[2]}
\UnaryInfC{$!B$}
\LeftSubproofLabel{$\delta_2$}
\AxiomC{$\Discharge{!A \land !B}{x}$}
\RightLabel{\Elim{\land}[1]}
\UnaryInfC{$!A$}
\RightSubproofLabel{$\delta_3$}
\RightLabel{\Intro{\land}}
\BinaryInfC{$!B \land !A$}
\DischargeRule{\Intro{\lif}}{x}
\UnaryInfC{$(!A \land !B) \lif (!B \land !A)$}
\DisplayProof
\]
तिची उंची
$\pheight{\delta_4} + 1 = 3$.
\Intro{\lif} अनुमान
$!A \land !B$
या रूपाची, खूण~$x$
असलेली सर्व गृहीतके
मुक्त करते; म्हणजे
निकटच्या उपसिद्धतेतील
आधी उघडी असलेली
दोन्ही गृहीतके.
~$\delta_5$ ची
तीच दोन गृहीतके असल्याने
तिच्यात उघडे गृहीतक
उरत नाही; म्हणून ती
तिच्या अंतिम-!!{formula}
$(!A \land !B)
\lif (!B \land !A)$
\emph{ची} !!a{proof}
आहे.
\end{ex}

\end{document}
